% Part: computability
% Chapter: computability-theory
% Section: ce-sets

\documentclass[../../../include/open-logic-section]{subfiles}

\begin{document}

\olfileid[pt]{cmp}{thy}{ces}
\olsection{Conjuntos computavelmente enumeráveis}

\begin{defn}
Um conjunto é \emph{computavelmente enumerável} se é vazio ou se é a
imagem de uma função computável.
\end{defn}

\begin{history}
Conjuntos computavelmente enumeráveis também são chamados de
\emph{recursivamente enumeráveis}. Essa é a terminologia original e,
hoje, tanto as duas denominações quanto as abreviações ``c.e.'' e
``r.e.'' são comumente usadas.
\end{history}

\begin{explain}
Você deve refletir sobre o significado da definição e sobre por que a
terminologia é apropriada. A ideia é que, se $S$ é a imagem da função
computável~$f$, então
\[
S = \{ f(0), f(1), f(2), \dots \},
\]
e, assim, $f$ pode ser vista como uma função que ``enumera'' os
elementos de~$S$. Observe que, segundo a definição, $f$~não precisa ser
uma função crescente, O.s.\@, a enumeração não precisa estar em ordem
crescente. Na verdade, $f$ nem sequer precisa ser injetora, O.s.\@,
são permitidas repetições na enumeração $f(0)$, $f(1)$, $f(2)$,
\dots{} de~$S$. Por exemplo, a função constante $f(x) = 0$ enumera o
conjunto $\{ 0 \}$.
\end{explain}

Todo conjunto computável é computavelmente enumerável. Para ver isso,
suponha que $S$~seja computável. Se $S$ for vazio, então, por definição,
ele será computavelmente enumerável. Caso contrário, seja $a$ um
elemento qualquer de $S$. Defina $f$ por
\[
f(x) =
\begin{cases}
x & \text{se $\Char{S}(x) = 1$} \\
a & \text{caso contrário.}
\end{cases}
\]
Então $f$ é uma função computável, e $S$ é a imagem de~$f$.

\end{document}
